% Generated draft: requires mathematical and source review.
\documentclass[11pt]{article}
\usepackage[margin=25mm]{geometry}
\usepackage{fontspec}
\setmainfont{Latin Modern Roman}
\usepackage{amsmath,amssymb,mathtools,unicode-math}
\setmathfont{Latin Modern Math}
\usepackage{enumitem,needspace,xurl,hyperref}
\setlength{\parindent}{0pt}
\setlength{\parskip}{5pt}
\setlength{\emergencystretch}{3em}
\newcommand{\N}{\mathbb{N}}
\newcommand{\Z}{\mathbb{Z}}
\newcommand{\Q}{\mathbb{Q}}
\newcommand{\R}{\mathbb{R}}
\newcommand{\C}{\mathbb{C}}
\newcommand{\F}{\mathbb{F}}
\newcommand{\A}{\mathbb{A}}
\newcommand{\PP}{\mathbb{P}}
\newcommand{\E}{\mathbb{E}}
\newcommand{\eps}{\varepsilon}
\newcommand{\dd}{\,\mathrm{d}}
\begin{document}
\section*{UnsolvedMath ID: 800659}
\textbf{Source problem code:} catalogue entry
\subsection{Definitions and mathematical statement}
For families \(\mathcal A,\mathcal B\) of subsets of a set \(E\), the selection principle \(\mathsf S_1(\mathcal A,\mathcal B)\) means that for every sequence \((A_n)_{n\geq 1}\) of members of \(\mathcal A\), there are choices \(a_n\in A_n\) such that \(\{a_n:n\geq 1\}\in\mathcal B\). The selection principle \(\mathsf S_{\mathrm{fin}}(\mathcal A,\mathcal B)\) means that for every such sequence there are finite sets \(F_n\subseteq A_n\) such that \(\bigcup_{n\geq 1}F_n\in\mathcal B\); the finite sets may be empty.

The supplied Problem 32 asks whether the properties
\[
\mathsf S_1(\mathcal D_\Omega,\mathcal D_\Omega)
\quad\text{or}\quad
\mathsf S_{\mathrm{fin}}(\mathcal D_\Omega,\mathcal D_\Omega)
\]
are preserved under taking finite powers. The record does not define \(\mathcal D_\Omega\), specify its underlying set of selectable objects, or identify the class of spaces under consideration. Consequently, a complete mathematical formulation cannot be recovered from the supplied record alone.

The preservation questions can be expressed schematically as follows. Let \(\mathscr C\) denote the intended class of spaces, and let \(\mathcal D_\Omega(X)\) denote the intended family associated with a space \(X\). For each \(\star\in\{1,\mathrm{fin}\}\), ask whether
\[
\forall X\in\mathscr C\;\forall m\in\mathbb Z_{\geq 1},\qquad
\mathsf S_\star\bigl(\mathcal D_\Omega(X),\mathcal D_\Omega(X)\bigr)
\Longrightarrow
\mathsf S_\star\bigl(\mathcal D_\Omega(X^m),\mathcal D_\Omega(X^m)\bigr).
\]
Here, if the intended setting is topological spaces, \(X^m\) denotes the Cartesian power equipped with the product topology. This is a conditional formulation template: the meanings of \(\mathscr C\) and \(\mathcal D_\Omega(X)\), including any restrictions on spaces or selectable families, remain unresolved. The two values of \(\star\) display the two properties mentioned in the source separately.
\subsection{Short English statement}
The source asks whether each of two selection properties involving \(\mathcal D_\Omega\) survives positive finite powers. Their selection rules are stated above, but the supplied record omits the definition of \(\mathcal D_\Omega\) and the intended class of spaces. Those omissions require clarification before this can be a complete mathematical statement.
\subsection{Sources}
\begin{itemize}
\item Local UnsolvedMath registry: permanent ID \(800659\), exact source problem code \texttt{catalogue entry}; title: \emph{Selection Principles Milestone Problems --- Problem 32 ([coc5])}. The supplied registry attributes the record to Tsaban, with proposed year 2003, and labels its status as open. This is a supplied status label, not an independently verified mathematical or literature assessment.
\item Supplied source citation: Tsaban, \emph{Selection Principles in Mathematics: A Milestone of Open Problems} (2003), Problem 32, \url{https://arxiv.org/abs/math/0312182}. The record reports access on 2026-10-16 and extraction by \texttt{arxiv\_latex}; these are supplied provenance details, not a claim of independent source access.
\item Supplied attribution/context label: \texttt{[coc5]}. Its bibliographic referent and any qualifications it supplies are absent from the record.
\end{itemize}
\end{document}
